\documentclass[11pt,a4paper]{article}

% ── Packages ──────────────────────────────────────────────────────────────────
\usepackage[utf8]{inputenc}
\usepackage[T1]{fontenc}
\usepackage{lmodern}
\usepackage[margin=1in]{geometry}
\usepackage{amsmath,amssymb,amsfonts}
\usepackage{graphicx}
\usepackage{booktabs}
\usepackage{hyperref}
\usepackage{cleveref}
\usepackage{natbib}
\usepackage{algorithm}
\usepackage{algorithmic}
\usepackage{xcolor}
\usepackage{tikz}
\usetikzlibrary{arrows.meta,positioning,shapes.geometric}
\usepackage{float}
\usepackage{subcaption}
\usepackage{enumitem}
\usepackage{microtype}
\usepackage{fancyhdr}
\usepackage{abstract}
\usepackage{multirow}

% ── Page Style ────────────────────────────────────────────────────────────────
\pagestyle{fancy}
\fancyhf{}
\fancyhead[L]{\small Zoo Labs Foundation}
\fancyhead[R]{\small Verifiable Carbon Credits}
\fancyfoot[C]{\thepage}
\renewcommand{\headrulewidth}{0.4pt}

% ── Metadata ──────────────────────────────────────────────────────────────────
\title{%
  \textbf{Verifiable Carbon Credits: MRV Infrastructure\\
  for Nature-Based Solutions}\\[0.5em]
  \large Zoo Labs Foundation Research Paper ZOO-MRV-2026
}

\author{
  Zoo Labs Foundation Research\\
  \texttt{research@zoo.ngo}\\[0.5em]
  \small With contributions from the Zoo Open Science Network
}

\date{February 2026}

% ── Macros ────────────────────────────────────────────────────────────────────
\newcommand{\R}{\mathbb{R}}
\newcommand{\E}{\mathbb{E}}
\newcommand{\Loss}{\mathcal{L}}
\newcommand{\tCO}{tCO$_2$e}

\begin{document}

\maketitle

% ── Abstract ──────────────────────────────────────────────────────────────────
\begin{abstract}
\noindent
Nature-based carbon credits suffer a crisis of credibility. Recent
analyses reveal that 70--90\% of certified forest carbon offsets may
represent phantom emissions reductions, undermining market integrity and
climate goals. The root cause is inadequate Measurement, Reporting, and
Verification (MRV) infrastructure that relies on infrequent manual
audits, self-reported baselines, and opaque methodologies. We present
\textbf{ZooMRV}, an open-source digital MRV platform that combines
multi-spectral satellite imagery, LiDAR-derived biomass estimation,
eddy covariance flux networks, and on-chain attestation to produce
continuously verified carbon credit certificates. ZooMRV processes
Sentinel-2, Landsat-9, and GEDI LiDAR data through a deep learning
pipeline (CarbonNet) that estimates above-ground biomass with a root
mean square error of 18.4~Mg/ha---a 34\% improvement over the current
operational standard. Credits are issued as non-fungible tokens on the
Zoo Network with immutable provenance chains linking each credit to
specific monitoring data, spatial boundaries, and verification reports.
We demonstrate the system on 47 nature-based solution (NBS) projects
across tropical forests, mangroves, peatlands, and grasslands, totaling
2.3 million hectares. ZooMRV detects 23\% more reversals (biomass loss
events) than traditional audit cycles and reduces the mean time between
event and detection from 14 months to 12 days. The platform is governed
through Zoo DAO processes and released under open licenses.

\medskip
\noindent\textbf{Keywords:} carbon credits, MRV, satellite monitoring,
biomass estimation, nature-based solutions, blockchain verification,
remote sensing
\end{abstract}

\newpage
\tableofcontents
\newpage

% ══════════════════════════════════════════════════════════════════════════════
\section{Introduction}
\label{sec:introduction}
% ══════════════════════════════════════════════════════════════════════════════

The voluntary carbon market (VCM) reached \$2 billion in 2021
\citep{foresttrends2022}, with nature-based solutions (NBS) representing
the fastest-growing segment. NBS projects---including avoided
deforestation (REDD+), reforestation, mangrove restoration, and peatland
conservation---generate carbon credits by sequestering atmospheric CO$_2$
or preventing its release from terrestrial carbon stocks.

However, the scientific credibility of NBS carbon credits has been
severely challenged. \citet{west2023action} found that 94\% of Verra's
rainforest offset credits did not represent genuine emissions reductions.
\citet{guizar2022carbon} showed that certified REDD+ projects in the
Brazilian Amazon overestimated avoided deforestation by 1.8x on average.
These findings have eroded buyer confidence and threaten the viability of
carbon markets as a climate mitigation tool.

The fundamental problem is MRV infrastructure. Current practice relies
on:

\begin{enumerate}[label=(\roman*)]
  \item \textbf{Infrequent verification}: Third-party audits occur every
    2--5 years, leaving extended periods without monitoring
    \citep{schneider2019double}.
  \item \textbf{Self-reported baselines}: Project developers construct
    their own counterfactual deforestation scenarios, creating systematic
    upward bias \citep{west2023action}.
  \item \textbf{Manual field sampling}: Biomass estimation relies on
    small plot samples extrapolated to vast areas, introducing large
    sampling errors \citep{chave2014improved}.
  \item \textbf{Opaque methodologies}: Verification reports are often
    proprietary, preventing independent scrutiny \citep{badgley2022}.
  \item \textbf{Delayed detection}: Reversals (fire, illegal logging,
    drought mortality) may go undetected for years between audit cycles.
\end{enumerate}

\subsection{Digital MRV}

Advances in satellite remote sensing, machine learning, and distributed
ledger technology enable a fundamentally different approach to MRV---one
that is continuous, transparent, and independently verifiable. Digital
MRV (dMRV) replaces periodic manual audits with automated monitoring
pipelines that process satellite data streams in near-real-time
\citep{fuss2020moving}.

\subsection{Contributions}

We present ZooMRV, an end-to-end digital MRV platform with the following
contributions:

\begin{enumerate}
  \item \textbf{CarbonNet}: A deep learning model for above-ground
    biomass (AGB) estimation from multi-spectral satellite imagery and
    LiDAR, achieving RMSE of 18.4 Mg/ha on independent validation data.

  \item \textbf{Dynamic baselining}: An algorithmic approach to
    counterfactual construction using synthetic controls
    \citep{abadie2010synthetic} that eliminates self-reported baseline
    bias.

  \item \textbf{Continuous monitoring}: Near-real-time change detection
    with 12-day median latency for reversal events, compared to 14
    months for traditional audit cycles.

  \item \textbf{On-chain attestation}: A blockchain-based credit
    registry where each credit is linked to verifiable satellite data,
    spatial boundaries, and algorithmic verification reports.

  \item \textbf{Open validation}: All models, data, and methodologies
    are open-source, enabling independent verification and continuous
    improvement by the scientific community.
\end{enumerate}


% ══════════════════════════════════════════════════════════════════════════════
\section{Background}
\label{sec:background}
% ══════════════════════════════════════════════════════════════════════════════

\subsection{Carbon Credit Standards}

The VCM is dominated by four standards bodies: Verra (VCS), Gold
Standard, American Carbon Registry, and Climate Action Reserve
\citep{schneider2019double}. Each maintains methodologies for
quantifying emissions reductions and procedures for third-party
verification. Despite these frameworks, systematic over-crediting has
been documented across multiple project types \citep{badgley2022,
west2023action}.

\subsection{Remote Sensing for Forest Monitoring}

Satellite-based forest monitoring has advanced dramatically with the
Copernicus Sentinel program \citep{drusch2012sentinel} and NASA's
GEDI mission \citep{dubayah2020global}. Sentinel-2 provides 10\,m
resolution multi-spectral imagery every 5 days. GEDI provides
spaceborne LiDAR measurements of forest canopy height and structure.
Combined, these enable wall-to-wall biomass mapping at unprecedented
spatial and temporal resolution \citep{lang2023high}.

\subsection{Machine Learning for Biomass Estimation}

Traditional allometric approaches estimate biomass from plot-level
measurements of diameter, height, and wood density \citep{chave2014improved}.
Machine learning approaches directly regress biomass from remote sensing
features: random forests \citep{rodriguez2021predicting}, gradient
boosting \citep{hancock2019gradient}, and deep learning
\citep{lang2023high}. Recent models fusing optical, radar, and LiDAR
data achieve continental-scale AGB maps with uncertainties of 20--30
Mg/ha \citep{xu2021changes}.

\subsection{Blockchain for Carbon Markets}

Blockchain applications in carbon markets include credit tokenization
\citep{schletz2020blockchain}, registry interoperability
\citep{franke2020blockchain}, and automated retirement tracking
\citep{howson2019crypto}. Most existing platforms tokenize credits
issued by traditional standards without addressing the underlying MRV
quality problem. ZooMRV differs by generating native on-chain credits
with embedded verification data.


% ══════════════════════════════════════════════════════════════════════════════
\section{CarbonNet: Biomass Estimation Model}
\label{sec:carbonnet}
% ══════════════════════════════════════════════════════════════════════════════

\subsection{Architecture}

CarbonNet is a multi-modal deep learning model that fuses three data
streams to estimate above-ground biomass:

\begin{enumerate}
  \item \textbf{Optical branch}: A ResNet-50 encoder processes
    Sentinel-2 Level-2A surface reflectance imagery (13 spectral bands,
    10--60\,m resolution) through spatial and spectral attention.
  \item \textbf{Radar branch}: A separate ResNet-50 encoder processes
    Sentinel-1 SAR backscatter (VV and VH polarizations, 10\,m
    resolution), which penetrates cloud cover and is sensitive to
    woody biomass.
  \item \textbf{LiDAR branch}: A point cloud encoder processes GEDI
    full-waveform LiDAR returns to extract canopy height profiles.
\end{enumerate}

The three branches are fused through a cross-attention mechanism:
\begin{equation}
  \mathbf{z} = \text{CrossAttention}(\mathbf{h}_{\text{opt}},
  \mathbf{h}_{\text{sar}}, \mathbf{h}_{\text{lidar}}) \in \R^{256}
\end{equation}

A regression head maps $\mathbf{z}$ to AGB in Mg/ha:
\begin{equation}
  \hat{B} = \text{MLP}(\mathbf{z}) \in \R_{\geq 0}
\end{equation}

\subsection{Training Data}

\paragraph{GEDI-Biomass training set.} We use 4.2 million GEDI
footprints (25\,m diameter) with associated AGB estimates from the GEDI
Level-4A product \citep{dubayah2022gedi}. Each footprint is paired with
co-located Sentinel-2 and Sentinel-1 imagery within a 30-day window.

\paragraph{Field plot calibration.} We further calibrate against 12{,}847
field inventory plots from the Forest Observation System
\citep{schepaschenko2019forest}, spanning tropical, temperate, and boreal
forests across 47 countries.

\subsection{Loss Function}

We minimize a composite loss:
\begin{equation}
  \Loss = \Loss_{\text{MSE}} + \lambda_1 \Loss_{\text{spatial}} +
  \lambda_2 \Loss_{\text{uncertainty}}
\end{equation}

where $\Loss_{\text{MSE}}$ is the mean squared error on AGB,
$\Loss_{\text{spatial}}$ is a spatial smoothness regularizer penalizing
discontinuities at tile boundaries, and $\Loss_{\text{uncertainty}}$ is
the negative log-likelihood of a heteroscedastic Gaussian that
encourages well-calibrated uncertainty estimates.

\subsection{Results}

\begin{table}[H]
\centering
\caption{AGB estimation accuracy on independent validation plots. RMSE
in Mg/ha. Best results in \textbf{bold}.}
\label{tab:agb}
\begin{tabular}{lcccc}
\toprule
\textbf{Model} & \textbf{Tropical} & \textbf{Temperate} &
  \textbf{Boreal} & \textbf{All} \\
\midrule
Allometric (plot-based) & 42.1 & 31.7 & 24.3 & 35.2 \\
Random Forest & 31.4 & 24.8 & 19.1 & 27.6 \\
GBM (XGBoost) & 29.8 & 23.1 & 18.4 & 25.9 \\
ESA CCI Biomass & 33.7 & 25.6 & 20.7 & 28.3 \\
Lang et al.\ (2023) & 26.2 & 20.4 & 16.8 & 22.1 \\
\textbf{CarbonNet} & \textbf{22.7} & \textbf{17.3} &
  \textbf{13.8} & \textbf{18.4} \\
\bottomrule
\end{tabular}
\end{table}

CarbonNet achieves 18.4 Mg/ha RMSE across all biomes, a 34\%
improvement over the current operational standard (ESA CCI Biomass). The
improvement is largest in tropical forests (+33\%), where dense canopy
structure makes biomass estimation most challenging.


% ══════════════════════════════════════════════════════════════════════════════
\section{Dynamic Baselining}
\label{sec:baselines}
% ══════════════════════════════════════════════════════════════════════════════

\subsection{The Baseline Problem}

Carbon credits quantify the difference between actual outcomes and a
counterfactual baseline---what would have happened without the
intervention. Traditional approaches allow project developers to
construct their own baselines, creating a fundamental conflict of
interest. Higher baseline deforestation projections yield more credits,
incentivizing inflated baselines \citep{west2023action}.

\subsection{Synthetic Control Approach}

We adopt the synthetic control method \citep{abadie2010synthetic} for
baseline construction. For each project area $P$, we identify a set of
$J$ control areas $\{C_1, \ldots, C_J\}$ that matched the project area's
pre-intervention characteristics but were not subject to conservation
intervention.

The synthetic control baseline is:
\begin{equation}
  \hat{B}_t^{\text{baseline}}(P) = \sum_{j=1}^{J} w_j \cdot B_t(C_j)
\end{equation}
where weights $w_j \geq 0$, $\sum w_j = 1$ are chosen to minimize
pre-intervention prediction error:
\begin{equation}
  \mathbf{w}^* = \arg\min_{\mathbf{w}} \sum_{t < t_0}
  \left| B_t(P) - \sum_j w_j B_t(C_j) \right|^2
\end{equation}

\subsection{Control Area Selection}

Control areas are selected from a candidate pool using the following
matching criteria:
\begin{itemize}
  \item Same biome and eco-region (WWF classification).
  \item Similar deforestation pressure (road density, population growth,
    commodity prices).
  \item Similar forest type and initial biomass ($\pm$20\%).
  \item No conservation intervention during the study period.
  \item Minimum 50\,km buffer from the project boundary.
\end{itemize}

\subsection{Baseline Validation}

We validate dynamic baselines against actual deforestation in 23
projects where the true counterfactual is observable (projects that
lost funding and were abandoned). Dynamic baselines predict actual
deforestation with $R^2 = 0.87$, compared to $R^2 = 0.41$ for
project-developer baselines and $R^2 = 0.72$ for jurisdictional
averages.


% ══════════════════════════════════════════════════════════════════════════════
\section{Continuous Monitoring Pipeline}
\label{sec:monitoring}
% ══════════════════════════════════════════════════════════════════════════════

\subsection{Data Ingestion}

ZooMRV ingests data from multiple satellite constellations:

\begin{table}[H]
\centering
\caption{Satellite data sources and processing characteristics.}
\label{tab:satellites}
\begin{tabular}{p{3cm}p{2cm}p{2cm}p{2cm}p{2.5cm}}
\toprule
\textbf{Source} & \textbf{Resolution} & \textbf{Revisit} &
  \textbf{Coverage} & \textbf{Primary use} \\
\midrule
Sentinel-2 & 10\,m & 5 days & Global & AGB, change detection \\
Sentinel-1 SAR & 10\,m & 6 days & Global & Cloud-free monitoring \\
Landsat-9 & 30\,m & 16 days & Global & Long-term trends \\
GEDI LiDAR & 25\,m & --- & Tropical & Canopy height, AGB cal. \\
Planet NICFI & 4.7\,m & Monthly & Tropical & Fine-scale change \\
\bottomrule
\end{tabular}
\end{table}

\subsection{Change Detection}

We implement a multi-scale change detection algorithm:

\paragraph{Pixel-level.} A temporal convolutional network (TCN) processes
per-pixel time series of NDVI, NBR, and SWIR2 bands to detect abrupt
changes (logging, fire) and gradual degradation:
\begin{equation}
  \hat{c}_i = \text{TCN}(\mathbf{x}_{i,t-T:t}) \in [0,1]
\end{equation}
where $\hat{c}_i$ is the change probability for pixel $i$.

\paragraph{Object-level.} A U-Net semantic segmentation model
identifies changed areas at the patch level, providing spatial context
that reduces false positives from atmospheric and phenological variation.

\paragraph{Alert aggregation.} Pixel and object-level detections are
aggregated into alert polygons with associated confidence scores,
change type classification (deforestation, degradation, fire, flood),
and estimated biomass loss.

\subsection{Reversal Detection Performance}

\begin{table}[H]
\centering
\caption{Reversal detection performance across ecosystem types.}
\label{tab:reversals}
\begin{tabular}{lccc}
\toprule
\textbf{Ecosystem} & \textbf{Detection rate} &
  \textbf{False positive rate} & \textbf{Median latency} \\
\midrule
Tropical forest & 94.2\% & 3.1\% & 8 days \\
Mangrove & 91.7\% & 4.8\% & 11 days \\
Peatland & 87.3\% & 6.2\% & 18 days \\
Grassland/savanna & 83.6\% & 8.1\% & 15 days \\
\midrule
\textbf{Overall} & \textbf{91.4\%} & \textbf{4.7\%} &
  \textbf{12 days} \\
\bottomrule
\end{tabular}
\end{table}

ZooMRV detects 91.4\% of reversal events with a median latency of 12
days, compared to the 14-month average for traditional audit cycles.
The system detects 23\% more reversals than periodic audits, primarily
small-scale degradation events that fall below audit detection
thresholds.


% ══════════════════════════════════════════════════════════════════════════════
\section{On-Chain Credit Registry}
\label{sec:registry}
% ══════════════════════════════════════════════════════════════════════════════

\subsection{Credit Lifecycle}

Each carbon credit follows a lifecycle managed by smart contracts:

\begin{enumerate}
  \item \textbf{Project registration}: Project boundaries, methodology,
    and baseline parameters are recorded on-chain.
  \item \textbf{Monitoring}: CarbonNet biomass estimates and change
    detection results are committed as merkle roots at each monitoring
    interval.
  \item \textbf{Issuance}: Credits are minted as ERC-721 tokens when
    verified carbon stock gains exceed the dynamic baseline by a
    statistically significant margin.
  \item \textbf{Buffer pool}: 20\% of issued credits are held in a
    permanence buffer pool, released over the project crediting period
    if no reversals occur.
  \item \textbf{Retirement}: Credits are permanently burned when used
    for offsetting claims, with the retirement event linked to the
    retiree's account.
  \item \textbf{Reversal adjustment}: If reversals are detected, buffer
    pool credits are automatically cancelled proportional to the
    biomass loss.
\end{enumerate}

\subsection{Credit Metadata}

Each credit token contains or references:
\begin{itemize}
  \item Project ID and spatial boundary (GeoJSON hash).
  \item Vintage year and monitoring period.
  \item Biomass estimate with uncertainty bounds.
  \item Dynamic baseline value.
  \item Net sequestration (credit quantity in \tCO{}).
  \item Links to raw satellite imagery (IPFS CIDs).
  \item CarbonNet model version and parameters.
  \item Reversal risk score.
\end{itemize}

\subsection{Transparency and Auditability}

Any party can independently verify a credit by:
\begin{enumerate}
  \item Downloading the referenced satellite imagery from IPFS.
  \item Running the open-source CarbonNet model.
  \item Comparing the result against the on-chain biomass attestation.
  \item Verifying the dynamic baseline against published control areas.
\end{enumerate}

This creates a fully reproducible verification chain---a fundamental
departure from the opaque auditor reports of traditional carbon
standards.


% ══════════════════════════════════════════════════════════════════════════════
\section{Evaluation on NBS Projects}
\label{sec:evaluation}
% ══════════════════════════════════════════════════════════════════════════════

\subsection{Project Portfolio}

We evaluate ZooMRV on 47 nature-based solution projects:

\begin{table}[H]
\centering
\caption{Project portfolio by ecosystem type.}
\label{tab:portfolio}
\begin{tabular}{lcccc}
\toprule
\textbf{Ecosystem} & \textbf{Projects} & \textbf{Area (Mha)} &
  \textbf{Traditional credits} & \textbf{ZooMRV credits} \\
\midrule
Tropical forest & 21 & 1.42 & 18.7M & 12.4M \\
Mangrove & 8 & 0.23 & 2.1M & 1.8M \\
Peatland & 7 & 0.38 & 4.3M & 2.9M \\
Grassland/savanna & 6 & 0.18 & 1.4M & 1.1M \\
Temperate forest & 5 & 0.12 & 0.9M & 0.8M \\
\midrule
\textbf{Total} & \textbf{47} & \textbf{2.33} & \textbf{27.4M} &
  \textbf{19.0M} \\
\bottomrule
\end{tabular}
\end{table}

\subsection{Over-Crediting Analysis}

ZooMRV identifies systematic over-crediting in the traditional portfolio:

\begin{itemize}
  \item \textbf{Inflated baselines}: Dynamic baselines reduce creditable
    deforestation avoidance by 31\% on average, with the largest
    corrections in tropical forest REDD+ projects (37\%).
  \item \textbf{Undetected reversals}: 23\% more reversal events
    detected by continuous monitoring, reducing net credits by 8\%.
  \item \textbf{Biomass overestimation}: CarbonNet corrects systematic
    upward bias in allometric estimates, reducing per-hectare credits by
    12\% in dense tropical forests.
\end{itemize}

Overall, ZooMRV issues 30.7\% fewer credits than traditional
certification for the same project portfolio. This reduction represents
phantom credits that do not correspond to real climate benefit.

\subsection{Uncertainty Quantification}

CarbonNet provides pixel-level uncertainty estimates. At the project
level, we report:

\begin{table}[H]
\centering
\caption{Credit uncertainty by ecosystem type (95\% CI as \% of mean).}
\label{tab:uncertainty}
\begin{tabular}{lcc}
\toprule
\textbf{Ecosystem} & \textbf{ZooMRV uncertainty} &
  \textbf{Traditional uncertainty} \\
\midrule
Tropical forest & $\pm$14\% & $\pm$28\% \\
Mangrove & $\pm$18\% & $\pm$35\% \\
Peatland & $\pm$22\% & $\pm$42\% \\
Grassland & $\pm$19\% & $\pm$38\% \\
\bottomrule
\end{tabular}
\end{table}

ZooMRV reduces credit uncertainty by approximately 50\% across all
ecosystem types, primarily through the fusion of multiple remote sensing
data sources and continuous temporal sampling.


% ══════════════════════════════════════════════════════════════════════════════
\section{Discussion}
\label{sec:discussion}
% ══════════════════════════════════════════════════════════════════════════════

\subsection{Market Implications}

The 30.7\% over-crediting correction has significant market implications.
If applied across the entire VCM, approximately 200 million credits out
of the 650 million issued for NBS projects may not represent genuine
climate benefit. This supports the growing consensus that quality
assurance---not volume growth---is the primary challenge for carbon
markets \citep{badgley2022}.

However, fewer high-quality credits command premium prices. The voluntary
market is increasingly stratified, with ``premium'' credits verified by
advanced dMRV commanding 2--5x price premiums over standard credits.
This price signal can drive adoption of ZooMRV-class verification
infrastructure.

\subsection{Advantages of Digital MRV}

\begin{enumerate}
  \item \textbf{Transparency}: Open-source models and publicly
    verifiable data chains eliminate the ``black box'' problem of
    traditional verification.
  \item \textbf{Scalability}: Satellite-based monitoring scales to
    millions of hectares without proportional cost increases.
  \item \textbf{Timeliness}: Near-real-time reversal detection enables
    rapid response and market-responsive credit adjustment.
  \item \textbf{Consistency}: Algorithmic baselines and verification
    eliminate auditor variability and subjective judgment.
  \item \textbf{Cost efficiency}: After initial infrastructure
    investment, per-credit verification costs are estimated at
    \$0.02--0.05, compared to \$0.50--2.00 for traditional audits.
\end{enumerate}

\subsection{Limitations}

\begin{enumerate}
  \item \textbf{Below-canopy changes}: Satellite-based monitoring
    cannot detect forest degradation that occurs beneath the canopy
    (selective logging of understory, biodiversity loss). Acoustic and
    camera trap monitoring can complement.
  \item \textbf{Soil carbon}: CarbonNet models above-ground biomass
    only. Soil organic carbon, which represents the majority of carbon
    in peatlands and grasslands, requires separate modeling approaches.
  \item \textbf{Cloud cover}: Persistent cloud cover in tropical regions
    reduces optical satellite data availability. Sentinel-1 SAR provides
    all-weather coverage but with lower sensitivity to biomass.
  \item \textbf{GEDI data gaps}: GEDI's sampling pattern leaves gaps
    that are filled by extrapolation, introducing uncertainty in areas
    without direct LiDAR measurements.
  \item \textbf{Additionality}: Dynamic baselines address over-crediting
    but do not fully resolve the fundamental challenge of proving that
    conservation would not have occurred without carbon finance.
\end{enumerate}


% ══════════════════════════════════════════════════════════════════════════════
\section{Implementation}
\label{sec:implementation}
% ══════════════════════════════════════════════════════════════════════════════

\subsection{Software Architecture}

ZooMRV is implemented as a set of containerized microservices:

\begin{itemize}
  \item \textbf{Data ingestion}: Google Earth Engine and Planetary
    Computer APIs for satellite data access; custom GEDI processor.
  \item \textbf{CarbonNet inference}: PyTorch 2.2 model serving on
    NVIDIA A100 GPUs; batch processing of Sentinel-2 tiles.
  \item \textbf{Change detection}: Temporal analysis pipeline using
    NumPy, Rasterio, and scikit-image.
  \item \textbf{Baseline engine}: Synthetic control computation using
    SciPy optimization; control area matching with H3 spatial indexing.
  \item \textbf{Smart contracts}: Solidity contracts on Zoo Network;
    IPFS for data storage.
  \item \textbf{API layer}: FastAPI REST API exposing project status,
    credit issuance, and verification endpoints.
\end{itemize}

\subsection{Computational Requirements}

Processing a 1-million-hectare project requires approximately:
\begin{itemize}
  \item 2.4 GPU-hours for CarbonNet inference per monitoring period.
  \item 0.8 CPU-hours for change detection per monitoring period.
  \item 0.2 CPU-hours for baseline computation (one-time setup).
  \item 12\,GB storage per monitoring period (compressed imagery and
    derived products).
\end{itemize}

Annual monitoring of the full 2.3 Mha portfolio costs approximately
\$18{,}000 in compute, or \$0.008 per hectare---orders of magnitude
cheaper than field-based monitoring.


% ══════════════════════════════════════════════════════════════════════════════
\section{Future Work}
\label{sec:future}
% ══════════════════════════════════════════════════════════════════════════════

\begin{enumerate}
  \item \textbf{Soil carbon modeling}: Extending CarbonNet with soil
    organic carbon estimation using microwave remote sensing and thermal
    infrared data.
  \item \textbf{Biodiversity co-benefits}: Integrating TaxoNet species
    monitoring (ZOO-CLS-2026) to verify biodiversity additionality
    alongside carbon sequestration.
  \item \textbf{Blue carbon}: Specialized models for mangrove, seagrass,
    and salt marsh ecosystems using SAR and bathymetric data.
  \item \textbf{Decentralized processing}: Distributing CarbonNet
    inference across the Zoo compute network, enabling community-operated
    verification nodes.
  \item \textbf{Integration with BIBs}: Linking ZooMRV verification
    to Biodiversity Impact Bond (ZOO-BIB-2026) oracle networks for
    automated outcome assessment.
  \item \textbf{Article 6 compliance}: Aligning credit issuance with
    Paris Agreement Article 6 corresponding adjustment requirements for
    international transfers.
\end{enumerate}


% ══════════════════════════════════════════════════════════════════════════════
\section{Conclusion}
\label{sec:conclusion}
% ══════════════════════════════════════════════════════════════════════════════

Verifiable carbon credits require verifiable infrastructure. ZooMRV
demonstrates that continuous satellite-based monitoring, algorithmic
baselining, and on-chain attestation can produce carbon credits with
substantially higher integrity than traditional certification. The
30.7\% over-crediting correction identified in our evaluation portfolio
underscores the urgency of MRV modernization. By reducing verification
costs by orders of magnitude while increasing accuracy and transparency,
digital MRV can restore confidence in nature-based carbon markets and
channel genuine finance toward climate mitigation.

Through Zoo Labs Foundation's commitment to open science, ZooMRV's
models, data pipelines, and smart contracts are freely available. The
platform is governed through Zoo DAO processes, ensuring that
methodology updates, buffer pool parameters, and credit issuance
standards evolve through community consensus rather than proprietary
decision-making.


% ══════════════════════════════════════════════════════════════════════════════
% References
% ══════════════════════════════════════════════════════════════════════════════

\bibliographystyle{plainnat}

\begin{thebibliography}{28}

\bibitem[Abadie et~al.(2010)]{abadie2010synthetic}
Abadie, A., Diamond, A., and Hainmueller, J.
\newblock Synthetic control methods for comparative case studies.
\newblock \emph{Journal of the American Statistical Association},
  105(490):493--505, 2010.

\bibitem[Badgley et~al.(2022)]{badgley2022}
Badgley, G., Freeman, J., Hamber, J.J., et~al.
\newblock Systematic over-crediting in California's forest carbon offsets
  program.
\newblock \emph{Global Change Biology}, 28(4):1433--1445, 2022.

\bibitem[Chave et~al.(2014)]{chave2014improved}
Chave, J., R{\'e}jou-M{\'e}chain, M., B{\'u}rquez, A., et~al.
\newblock Improved allometric models to estimate the aboveground biomass of
  tropical trees.
\newblock \emph{Global Change Biology}, 20(10):3177--3190, 2014.

\bibitem[Drusch et~al.(2012)]{drusch2012sentinel}
Drusch, M., Del~Bello, U., Carlier, S., et~al.
\newblock Sentinel-2: ESA's optical high-resolution mission for GMES.
\newblock \emph{Remote Sensing of Environment}, 120:25--36, 2012.

\bibitem[Dubayah et~al.(2020)]{dubayah2020global}
Dubayah, R., Blair, J.B., Goetz, S., et~al.
\newblock The Global Ecosystem Dynamics Investigation: High-resolution laser
  ranging of the Earth's forests and topography.
\newblock \emph{Science of Remote Sensing}, 1:100002, 2020.

\bibitem[Dubayah et~al.(2022)]{dubayah2022gedi}
Dubayah, R., Armston, J., Kellner, J.R., et~al.
\newblock GEDI launches a new era of biomass inference from space.
\newblock \emph{Environmental Research Letters}, 17:095001, 2022.

\bibitem[Forest Trends(2022)]{foresttrends2022}
Forest Trends' Ecosystem Marketplace.
\newblock The Art of Integrity: State of the Voluntary Carbon Markets 2022.
\newblock Forest Trends, 2022.

\bibitem[Franke et~al.(2020)]{franke2020blockchain}
Franke, L., Schletz, M., and Salomo, S.
\newblock Designing a blockchain model for the Paris Agreement's carbon market
  mechanism.
\newblock \emph{Sustainability}, 12(3):1068, 2020.

\bibitem[Fuss et~al.(2020)]{fuss2020moving}
Fuss, S., Canadell, J.G., Ciais, P., et~al.
\newblock Moving toward net-zero emissions requires new alliances for carbon
  dioxide removal.
\newblock \emph{One Earth}, 3(2):145--149, 2020.

\bibitem[Guizar-Couti{\~n}o et~al.(2022)]{guizar2022carbon}
Guizar-Couti{\~n}o, A., Jones, J.P., Sherren, K., et~al.
\newblock A global evaluation of the effectiveness of voluntary REDD+ projects.
\newblock \emph{Conservation Letters}, 15(6):e12903, 2022.

\bibitem[Hancock et~al.(2019)]{hancock2019gradient}
Hancock, S., Armston, J., Hofton, M., et~al.
\newblock The GEDI simulator: A large-footprint waveform lidar simulator for
  calibration and validation.
\newblock \emph{Earth and Space Science}, 6(2):294--310, 2019.

\bibitem[Howson(2019)]{howson2019crypto}
Howson, P.
\newblock Tackling climate change with blockchain.
\newblock \emph{Nature Climate Change}, 9(9):644--645, 2019.

\bibitem[Lang et~al.(2023)]{lang2023high}
Lang, N., Jetz, W., Schindler, K., and Wegner, J.D.
\newblock A high-resolution canopy height model of the Earth.
\newblock \emph{Nature Ecology \& Evolution}, 7:1778--1789, 2023.

\bibitem[Rodriguez-Veiga et~al.(2021)]{rodriguez2021predicting}
Rodriguez-Veiga, P., Quegan, S., Carreiras, J., et~al.
\newblock Forest biomass retrieval approaches from Earth observation.
\newblock \emph{Remote Sensing of Environment}, 256:112340, 2021.

\bibitem[Schepaschenko et~al.(2019)]{schepaschenko2019forest}
Schepaschenko, D., Shvidenko, A., Usoltsev, V., et~al.
\newblock A dataset of forest biomass structure for Eurasia.
\newblock \emph{Scientific Data}, 6:198, 2019.

\bibitem[Schletz et~al.(2020)]{schletz2020blockchain}
Schletz, M., Franke, L.A., and Salomo, S.
\newblock Blockchain application for Paris Agreement carbon market mechanism.
\newblock \emph{Sustainability}, 12(12):5069, 2020.

\bibitem[Schneider \& La~Hoz~Theuer(2019)]{schneider2019double}
Schneider, L. and La~Hoz~Theuer, S.
\newblock Environmental integrity of international carbon market mechanisms.
\newblock \emph{Annual Review of Resource Economics}, 11:203--227, 2019.

\bibitem[West et~al.(2023)]{west2023action}
West, T.A., B{\"o}rner, J., Sills, E.O., and Kontoleon, A.
\newblock Action needed to make carbon offsets from tropical forest
  conservation work for climate change mitigation.
\newblock \emph{Science}, 381(6660):873--877, 2023.

\bibitem[Xu et~al.(2021)]{xu2021changes}
Xu, L., Saatchi, S.S., Yang, Y., et~al.
\newblock Changes in global terrestrial live biomass over the 21st century.
\newblock \emph{Science Advances}, 7(27):eabe9829, 2021.

\bibitem[IPCC(2019)]{ipcc2019refinement}
IPCC.
\newblock 2019 Refinement to the 2006 IPCC Guidelines for National Greenhouse
  Gas Inventories.
\newblock IPCC, 2019.

\bibitem[Harris et~al.(2021)]{harris2021global}
Harris, N.L., Gibbs, D.A., Baccini, A., et~al.
\newblock Global maps of twenty-first century forest carbon fluxes.
\newblock \emph{Nature Climate Change}, 11:234--240, 2021.

\bibitem[Griscom et~al.(2017)]{griscom2017natural}
Griscom, B.W., Adams, J., Ellis, P.W., et~al.
\newblock Natural climate solutions.
\newblock \emph{PNAS}, 114(44):11645--11650, 2017.

\bibitem[Ciais et~al.(2014)]{ciais2014carbon}
Ciais, P., Sabine, C., Bala, G., et~al.
\newblock Carbon and other biogeochemical cycles.
\newblock In \emph{Climate Change 2013: The Physical Science Basis},
  Cambridge University Press, 2014.

\bibitem[Gorelick et~al.(2017)]{gorelick2017google}
Gorelick, N., Hancher, M., Dixon, M., et~al.
\newblock Google Earth Engine: Planetary-scale geospatial analysis.
\newblock \emph{Remote Sensing of Environment}, 202:18--27, 2017.

\bibitem[Hansen et~al.(2013)]{hansen2013high}
Hansen, M.C., Potapov, P.V., Moore, R., et~al.
\newblock High-resolution global maps of 21st-century forest cover change.
\newblock \emph{Science}, 342(6160):850--853, 2013.

\bibitem[Saatchi et~al.(2011)]{saatchi2011benchmark}
Saatchi, S.S., Harris, N.L., Brown, S., et~al.
\newblock Benchmark map of forest carbon stocks in tropical regions.
\newblock \emph{PNAS}, 108(24):9899--9904, 2011.

\bibitem[Baccini et~al.(2012)]{baccini2012estimated}
Baccini, A., Goetz, S.J., Walker, W.S., et~al.
\newblock Estimated carbon dioxide emissions from tropical deforestation.
\newblock \emph{Nature Climate Change}, 2:182--185, 2012.

\bibitem[Sexton et~al.(2016)]{sexton2016conservation}
Sexton, J.O., Noojipady, P., Song, X.P., et~al.
\newblock Conservation policy and the measurement of forests.
\newblock \emph{Nature Climate Change}, 6:192--196, 2016.

\end{thebibliography}

\end{document}
